\chapter{Projektüberblick}
\section{Zielsetzung}
	Da die Entwicklung unseres 6dof Roboterarms im Rahmen des Faches ASIT erfolgt und der Lernerfolg im Vordergrund steht konnten wir uns die Aufgabenstellung selbst stellen und hatten dabei keine großen Einschränkungen.  Wir haben uns dazu entschieden das der Roboterarm als Mindestanforderung eine Pick and Place Operation, also ein Objekt aufheben und wo anders platzieren können soll. Als optionale weitere Herausforderung soll man über ein Graphical user Interface Koordinaten eingeben können zu denen sich der Arm dann bewegt. Nach einem Schuljahr Arbeit an dem Projekt haben wir uns entschieden es als Diplom und Abschlussarbeit fortzusetztn.

\section{Umsetzung}
\section{Umsetzungsstand der Software}
\label{sec:umsetzungsstand}
Der aktuelle Umsetzungsstand der Software umfasst ein vollständiges Initialisieren der Motortreiber und ein sequentielles Homing. Dieses sequentielle Homing wird im verlauf des Projekts aber noch durch ein Absolut encoderhoming ersetzt.

%\chapter{Einleitung}